\chapter{KAJIAN ILMIAH}

\section{Kubernetes}

Kubernetes adalah platform orkestrasi kontainer \textit{open-source} yang dirancang untuk mengotomatisasi deployment, penskalaan, dan pengelolaan aplikasi kontainer \citep{kubernetes2022}. Kubernetes menyediakan abstraksi yang kuat untuk mengelola kompleksitas infrastruktur modern.

\subsection{Arsitektur Kubernetes}

Arsitektur Kubernetes terdiri dari Control Plane dan Worker Nodes. Control Plane mengelola status klaster melalui komponen seperti API Server, etcd, Scheduler, dan Controller Manager. Worker Nodes menjalankan aplikasi kontainer menggunakan kubelet dan kube-proxy.

\subsection{Manifest Kubernetes}

Semua objek Kubernetes didefinisikan secara deklaratif menggunakan manifest YAML. Pendekatan deklaratif ini menjadi fondasi bagi praktik GitOps \citep{kubernetes-docs}.

\section{GitOps}

GitOps adalah paradigma operasional yang menggunakan Git sebagai sumber kebenaran tunggal untuk infrastruktur dan aplikasi \citep{weaveworks-gitops}.

\subsection{Prinsip GitOps}

Menurut OpenGitOps Working Group, terdapat empat prinsip utama GitOps \citep{gitops-principles}:

\begin{enumerate}
    \item \textbf{Declarative:} Sistem dideskripsikan secara deklaratif.
    \item \textbf{Versioned and Immutable:} Status yang diinginkan disimpan dalam versi yang tidak dapat diubah.
    \item \textbf{Pulled Automatically:} Software agent secara otomatis menarik perubahan dari sumber kebenaran.
    \item \textbf{Continuously Reconciled:} Software agent terus-menerus memantau dan mengoreksi penyimpangan.
\end{enumerate}

\section{ArgoCD}

ArgoCD adalah \textit{GitOps controller} \textit{open-source} untuk Kubernetes yang mengotomatiskan deployment aplikasi \citep{argocd-docs}.

\section{Tinjauan Keamanan}

Keamanan dalam implementasi GitOps menjadi perhatian penting \citep{gitops-security}. Aspek keamanan yang relevan meliputi secrets management, network policies, dan pod security.

\section{InvenioRDM}

InvenioRDM adalah platform manajemen data riset \textit{open-source} yang dibangun di atas framework Invenio.
